%% kapa: exam=1 round=21 year=2010 period=2 form=A subject=appraisal_law subject_ko=부동산관계법규 seq=24 number=64
%% source: Q-Net 감정평가사 기출문제 (cst00309) artlSeq=5024181 file=extracted/1008/5024181/2010년 제21회 감정평가사 1차시험 기출문제/2교시.hwp
%% conversion: kordoc 4.24.51 -> tools/gen1.py (auto); review=proofread-gpt-sol; flags=none
\begin{question}{64}
부동산 가격공시 및 감정평가에 관한 법령상 개별공시지가에 대한 불복에 관한 설명으로 옳은 것은? (단, 다툼이 있으면 판례에 의함)
\begin{choices}
  \item 개별공시지가결정에 대하여 곧바로 행정소송을 제기할 수 없다.
  \item 개별공시지가의 결정․공시일부터 30일 이내에 서면으로 이의신청할 수 있다.
  \item 시장․군수 또는 구청장은 이의신청을 심사하기 위하여 반드시 감정평가업자에게 검증을 의뢰하여야 한다.
  \item 시장․군수 또는 구청장은 이의신청에 대한 심사 결과를 신청인에게 구두로 통지할 수 있다.
  \item 개별공시지가결정은 취소소송의 대상으로서의 처분성이 인정되지 아니한다.
\end{choices}
\end{question}
\answer{2}
